\documentclass[10pt]{scrartcl}
\usepackage[T1]{fontenc}	% Passendes Fontencoding zur Suche im pdf
\usepackage{textcomp}		% das EUR-Zeichen für OT und T1
\usepackage{lmodern}
\usepackage[ngerman]{babel}	% neue deutsche Rechtschreibung
\usepackage[utf8]{inputenc}	% direkt deutsche Umlaute und EUR-Zeichen eingeben
\usepackage{graphicx,framed}	% Bilder, Rahmen
\usepackage{wrapfig}
\usepackage[a6paper,paperheight=21cm]{geometry}
\usepackage{fancyhdr}
\usepackage{url}
\usepackage{tikz}
\usepackage{draftwatermark}

\newcommand{\ffmagenta}{magenta}
\newcommand{\ffyellow}{yellow}
\newcommand{\fflogo}{
\draw (36mm,44mm) circle (22mm) [line width=1mm,color=\ffmagenta];
\draw (74mm,54mm) circle (22mm) [line width=1mm,color=\ffmagenta];
\draw (74mm,54mm) circle (29mm) [line width=6mm,color=\ffmagenta];
\fill[\ffyellow] (22mm,48mm) -- (22mm,52mm) -- (30mm,52mm) -- (26mm,56mm) -- (31mm,56mm) -- (37mm,50mm) -- (31mm,44mm) -- (26mm,44mm) -- (30mm,48mm);
\fill[\ffyellow] (64mm,50mm) -- (73mm,59mm) -- (81mm,59mm) -- (75mm,53mm) -- (88mm,53mm) -- (88mm,47mm) -- (75mm,47mm) -- (81mm,41mm) -- (73mm,41mm);
}

\SetWatermarkAngle{0}
\renewcommand{\ffmagenta}{magenta!25!white}
\renewcommand{\ffyellow}{yellow!25!white}
\SetWatermarkText{
\begin{tikzpicture}[scale=1.7,rotate=90]
\fflogo
\end{tikzpicture}
}

\pagestyle{fancy}
\fancyhead{
\begin{tikzpicture}
\draw (\topmargin,0) -- (\textwidth,0) [line width=5mm,color=yellow];
\draw (0,-0.5) -- (\textwidth,-0.5) [line width=5mm,color=magenta];
\node at (5.5,0) {freifunk rheinland e.V.};
\node at (6,-0.5) {funkzelle aachen};
\end{tikzpicture}
}
\fancyfoot{}
\renewcommand{\headrule}{}
\renewcommand{\footrule}{}

\setcounter{secnumdepth}{0}

\begin{document}
\section{\normalsize Was ist Freifunk?}
Freifunk ermöglicht jedem Bürger mittels einfacher WLAN-Technik einen 
kostenfreien, unabhängigen Zugang zum Internet.
Dieses Projekt wird ehrenamtlich aufgebaut und betreut.

Das Ziel ist, ein stadtweit übergreifendes, nichtkommerzielles, für alle frei 
zugängiges Funknetz aufzubauen.
Jeder kann sich dem anschließen, es ist kein Technikwissen erforderlich.

\section{\normalsize Warum sollte ich freifunken?}
Ist Dir bewusst, wie oft Du Dein Internet nicht benutzt? Verbessere einfach die Erreichbarkeit des Freifunknetzes oder teile Dein Internet und lasse damit sozial benachteiligte Menschen, Besucher und frisch Zugezogene teilhaben und ermögliche ihnen Kommunikation und Zugang zu Wissen.

\section{\normalsize Ein Verein für Freifunk?}
\begin{wrapfigure}{R}{70px}
\includegraphics[scale=0.2]{Verein}
\end{wrapfigure}
Freifunker aus vielen Städten in NRW haben sich zusammengeschlossen und diesen Verein ins Leben gerufen, um Freifunk in den Regionen des Rheinlands zu fördern und bekannter zu machen.
\textbf{Hilf uns dabei!}
\lfoot{\tiny v.i.S.d.P. freifunk rheinland e.V., Postfach 27 01 02, 40524 Düsseldorf}

\newpage
\lfoot{\footnotesize\url{www.freifunk-rheinland.net}\\\url{kontakt@freifunk-rheinland.net}\\\tiny\centering Konto 4 044 542 600, BLZ 430 609 67, GLS Gemeinschaftsbank eG}
\rfoot{\footnotesize\url{@ffrhein}\\02 11/41 74 12 21}

\section{\normalsize Vorteile unseres Vereins}
Unter dem Dach eines gemeinnützigen Vereins fällt es uns leichter, an die Öffentlichkeit zu treten und Unterstützer zu finden. Wir bieten unseren Mitgliedern zudem Zugang zu Know-How, Infrastruktur und rechtlichen Schutz.

\section{\normalsize Du kannst uns helfen!}
\begin{wrapfigure}{R}{60px}
\includegraphics[scale=0.4]{Router}
\end{wrapfigure}
Stelle einen Freifunkrouter bei Dir auf, den Du von uns bekommen kannst. Wir sind aber auch auf Spenden und Mitgliedsbeiträge angewiesen, um die nötigen WLAN-Geräte und Zubehör besorgen zu können.

Hilf uns, neue Standorte für den Netzausbau zu finden. Ideal ist, wenn sie hoch und zentral platziert sind. Kontakte zu Hausverwaltern/-besitzern, Kirchen etc. sind sehr hilfreich.

Erzähle anderen von unserem Projekt!

\section{\normalsize Du möchtest uns persönlich kennenlernen?}
Komm einfach zu einem unserer Treffen oder nimm Kontakt zu uns auf:
\begin{framed}
\noindent An jedem ersten Dienstag des Monats um 19:30h in der Jülicher Str. 191 in Aachen, im Hackerspace des CCCAC.
\end{framed}
\end{document}
